\section{Preliminaries}
\label{sec:prelim}

\subsection{Koblitz curves and the Frobenius endomorphism}

Let $n$ be an odd prime and $K=\Fq$. The two Koblitz (anomalous binary)
curves are
\[
  E_a:\; y^2+xy = x^3+ax^2+1,\qquad a\in\{0,1\},
\]
defined over $\FF_2$ \citep{koblitz1992}. The $2$-power Frobenius
$\tau(x,y)=(x^2,y^2)$ is an endomorphism of $E_a$ satisfying
\[
  \tau^2-\mu\tau+2=0,\qquad \mu=(-1)^{1-a},
\]
and $\tau^n$ is the identity on $E_a(K)$. Negation is
$-(x,y)=(x,\,x+y)$, so the only points with $-P=P$ are $O$ and the
$2$-torsion point $(0,1)$. Throughout, $G\le E_a(K)$ denotes a cyclic
subgroup of prime order $r$ with $r>4$; for the fields in this paper the
cofactor $\#E_a(K)/r$ is $4$ (for $a=0$) or $2$ (for $a=1$).

Because $\tau$ commutes with scalar multiplication it preserves $G$, and
because $G$ is cyclic of prime order it acts on $G$ as a scalar:
there is a unique $\lambda\in(\ZZ/r\ZZ)^\times$ with
$\tau(P)=[\lambda]P$ for all $P\in G$. Reducing the characteristic
equation, $\lambda^2-\mu\lambda+2\equiv 0 \pmod r$ and
$\lambda^n\equiv 1\pmod r$. The scalar $\lambda$ is the quantity that
lets a Frobenius image be folded into a logarithm:
$\log_P(\tau^k Q)=\lambda^k\log_P Q$.

\subsection{The signed-Frobenius group}

Let $\Gamma=\langle\tau,-1\rangle$ be the group generated by $\tau$ and
negation acting on $E_a(K)$. Since $\tau^n=1$, $(-1)^2=1$ and the two
commute, $\Gamma$ is a quotient of $\ZZ/n\ZZ\times\ZZ/2\ZZ$ and has
order at most $2n$. Its orbits on $G\setminus\{O\}$ are the
\emph{signed-Frobenius orbits}; a factor base closed under $\Gamma$ is
\emph{$\Gamma$-invariant}. The following standard fact is the
foundation of everything that follows, and we state it with its
hypothesis made explicit because the ``exceptional short orbit'' case is
exactly the case in which it fails.

\begin{lemma}[Free action]
\label{lem:free}
Suppose that $r\nmid\#E_a(\FF_{2^d})$ for every proper divisor $d$ of
$n$. Then $\Gamma$ acts freely on $G\setminus\{O\}$: if
$\pm\tau^k(P)=P$ for some $P\in G\setminus\{O\}$ then $n\mid k$ and the
sign is $+$. In particular every signed-Frobenius orbit in
$G\setminus\{O\}$ has exactly $2n$ elements, and the
$x$-coordinates $x,x^2,\dots,x^{2^{n-1}}$ of its members are pairwise
distinct. When $n$ is prime the hypothesis reduces to $r>4$.
\end{lemma}

\begin{proof}
Let $d=\ord_r(\lambda)$. Since $\lambda^n\equiv 1$, $d\mid n$. If
$d<n$ then $\tau^d$ fixes every point of $G$, and the fixed points of
$\tau^d$ on $E_a(\bar K)$ are exactly $E_a(\FF_{2^d})$, so
$G\subseteq E_a(\FF_{2^d})$ and $r\mid\#E_a(\FF_{2^d})$, contradicting
the hypothesis. Hence $d=n$, which is odd. Now
$\tau^k(P)=\pm P$ gives $\lambda^k\equiv\pm1$, so
$\lambda^{2k}\equiv 1$, so $n\mid 2k$, so $n\mid k$ since $n$ is odd,
and then $\lambda^k\equiv 1\ne -1$ because $r$ is odd. The claim about
$x$-coordinates follows: $x^{2^i}=x^{2^j}$ with $0\le i<j<n$ would give
$\tau^{j-i}(P)\in\{P,-P\}$. For $n$ prime the only proper divisor is
$d=1$ and $\#E_a(\FF_2)\in\{2,4\}$.
\end{proof}

\begin{remark}[Exceptional short orbits]
\label{rem:short}
The points that do have short orbits are exactly those excluded by the
lemma: $O$; the $2$-torsion point $(0,1)$, fixed by both $\tau$ and
$-1$; the points of $E_a(\FF_2)$ (which for $a=0$ are
$(1,0),(1,1)$, interchanged by negation and fixed by $\tau$); and, for
composite $n$, the points of the proper subfield curves
$E_a(\FF_{2^d})$. None of them lies in $G$ when the hypothesis holds,
but all of them do arise as $x$-coordinates in a decomposition
pipeline (a summation-polynomial root $x\in\FF_2$ is not rare at small
$n$, and $x=0$ corresponds to a $2$-torsion ``solution''). A
canonicalization routine therefore has to be \emph{total} on
$\FF_{2^n}$ and must flag these words rather than mis-assign a phase to
them. Section~\ref{sec:canon-correctness} makes this precise.
\end{remark}

\subsection{Normal bases and rotation}

Fix a normal basis $\{\beta,\beta^2,\dots,\beta^{2^{n-1}}\}$ of $K$
over $\FF_2$ and write $x=\sum_i w_i\beta^{2^i}$ with
$w=(w_0,\dots,w_{n-1})\in\FF_2^n$. Then $x^2$ has coordinate vector
$\rot_1(w)=(w_{n-1},w_0,\dots,w_{n-2})$: squaring is a cyclic
rotation, and the Frobenius orbit of $x$ is the \emph{necklace} of
$w$. The element $1$ has coordinate vector $1^n$ and $0$ has $0^n$;
these are the only words fixed by every rotation, and for $n$ prime
they are the only words with any nontrivial rotational symmetry
(a word with $\rot_j(w)=w$, $0<j<n$, has period $\gcd(j,n)=1$). This is
the normal-basis form of Lemma~\ref{lem:free}: for prime $n$ every
word outside $\{0^n,1^n\}$ has exactly $n$ distinct rotations.

For the field sizes in this paper the normal basis is a Gaussian
normal basis of type~$2$ (an optimal normal basis, ONB) when one
exists ($n=131$ for ECC2K-130, $n=83$); the repository's packed
arithmetic computes multiplication in a polynomial basis and converts,
so that squaring, Frobenius powers and the canonicalization below cost
only word rotations~\citep{bailey2009}. A permutation basis-change costs
one table pass per direction and is charged in every pipeline count
reported here.

\subsection{Index calculus with a factor base of orbits}
\label{sec:ic-recap}

Let $F\subset G$ be a factor base and $P,Q\in G$ with $Q=[s]P$ to be
found. Relation collection draws random $R=[\alpha]P+[\beta]Q$ and
tries to decompose $R=\sum_{i=1}^m\pm P_i$ with $P_i\in F$ by
Semaev's summation polynomial $S_{m+1}$ \citep{semaev2004}, restricted
to a subspace of $x$-coordinates \citep{gaudry2009,diem2011}. Each
success is a linear relation among the unknowns
$\{\log_P P_i\}$, and once the relation matrix has rank $|F|$ (after
cofactor clearing) a solve modulo $r$ yields the logarithms and, from a
relation involving $Q$, the answer $s$.

If $F$ is $\Gamma$-invariant it is a union of $|F|/2n$ orbits. Choose a
representative $\hat P_c$ in each orbit $c$. Every $P_i\in F$ is
$P_i=\epsilon_i\tau^{k_i}(\hat P_{c(i)})$ for a unique sign
$\epsilon_i\in\{\pm1\}$ and phase $k_i\in\ZZ/n\ZZ$, so
\[
  \log_P P_i=\epsilon_i\lambda^{k_i}\log_P\hat P_{c(i)},
\]
and a decomposition becomes a relation among the $|F|/2n$ unknowns
$\log_P\hat P_c$ with coefficients in the order-$2n$ subgroup
$\langle\lambda,-1\rangle\le(\ZZ/r\ZZ)^\times$. The number of columns
falls by $2n$ while the relation yield per attempt is unchanged, so the
number of relations needed, and hence the number of decomposition
attempts, falls by $2n$ and the linear-algebra cost by roughly $(2n)^2$.
This is the factor-base symmetry exploited by
\citet{wienerzuccherato1998,glv2001} for Pollard rho and by
\citet{gg2014,fghr2014} for index calculus; we
recall it only to fix the notation and to make clear that the paper
does not claim it. What the paper is about is the \emph{cost of
realizing it}: finding the orbit representative and the phase of a
decomposed point, storing representatives compactly, probing them, and
carrying the phase through relation reconstruction and verification
without losing it.
